% Lamb_Sumit.tex --- master file, complete hand calculation.
% Compile: pdflatex, bibtex, pdflatex, pdflatex.
% \showworkedtrue  -> master with the working shown
% \showworkedfalse -> scholar version, results and sources only
\documentclass[11pt]{article}
\usepackage[T1]{fontenc}
\usepackage{lmodern}
\usepackage{amsmath,amssymb}
\usepackage{amsthm}
\usepackage{booktabs}
\usepackage{array}
\usepackage[a4paper,margin=2.2cm]{geometry}
\usepackage[hidelinks]{hyperref}

\newif\ifshowworked
\showworkedtrue
\makeatletter
\newenvironment{worked}{%
  \ifshowworked\par\smallskip\noindent\textit{Working.}\par\smallskip\else
  \setbox\z@\vbox\bgroup\fi}{\ifshowworked\par\smallskip\else\egroup\fi}
\makeatother

\newcommand{\ii}{\mathrm{i}}
\newcommand{\dd}{\mathrm{d}}
\newcommand{\Ncal}{\mathcal{N}}
\newcommand{\Gcal}{\mathcal{G}}
\setlength{\parskip}{3pt}
\setlength{\parindent}{0pt}

\title{\vspace{-1.2cm}Guided Lamb waves in a piezoelectric semiconductor plate:\\
carrier transport, nonlocal strain gradient elasticity and a heat equation with one lag}
\author{Dr.\ Vipin \\ Gurugram University}
\date{30 October 2026}

\begin{document}
\maketitle

\begin{abstract}
\noindent
We work out the plane-strain Lamb problem for a plate of thickness $2h$ cut from a piezoelectric
semiconductor of the ZnO family \cite{Ozgur2005}, with free faces at $z=\pm h$ and the six-fold axis along the thickness.
Five fields are carried: the two displacements $u_x,u_z$, the electric potential $\varphi$, the excess
carrier density $\Delta n$ and the temperature rise $\theta$. The mechanical stress follows the nonlocal
strain gradient law, with one length $l_{\rm nl}$ softening and one length $l_g$ stiffening, so the size
effect enters only as a length over the thickness. The heat equation is of the Moore--Gibson--Thompson
type in its base form, one phase lag and one time derivative cancelled. The carrier equation carries
drift, diffusion, thermodiffusion and an optional recombination term, and the heat equation carries the
reciprocal Peltier term. The five field equations are reduced to a $5\times5$ matrix in the thickness
exponent, its characteristic polynomial is counted and its parity decided, the fourteen partial waves are
written with their amplitude ratios, the seven face conditions are derived, and the dispersion matrix is
assembled. Two hand limits close the file.
\end{abstract}

\noindent\textbf{Keywords:} guided waves; piezoelectric semiconductor; nonlocal strain gradient elasticity;
thermoelastic coupling; carrier transport; attenuation.

\section{Introduction}

A Lamb wave is a guided elastic disturbance in a plate: the two free faces reflect it back and forth, and
the plate carries it along its own plane with a wavelength and a decay that depend on the product of
frequency and thickness. The plate is the natural element of acoustic filters, delay lines and sensors,
and in a piezoelectric material the wave carries an electric potential with it, so the same element both
senses and drives.

A piezoelectric \emph{semiconductor} differs from a piezoelectric dielectric in one respect that matters
here. The electric field is not free to arrange itself, because mobile carriers drift and diffuse under it
and screen it. Yang and Zhou wrote the drift--diffusion pair for a piezoelectric semiconductor plate and
used it on acoustic amplification
\cite[Eqs.~(3), (7)]{YangZhou2005}. The equations themselves are classical semiconductor physics; what
they add to the plate problem is a fourth field with its own length scale, the screening length. Guided
waves in semiconducting plates of graded or layered cross-section were taken up later with numerical
machinery that carries a variable coefficient through the thickness \cite{Othmani2017}.

The thermal side is partly occupied too. Sharma, Pal and Chand solved the piezothermoelastic plate with
free, electrically shorted and thermally insulated faces and got closed secular equations for the two
families of modes \cite{Sharma2004}. Kumar and Kansal added a diffusing species to a transversely
isotropic thermoelastic plate \cite{KumarKansal2008}. In both, the third field is a concentration that
carries no charge, so the electric and the diffusive problems do not talk to each other as they do in a
semiconductor.

The two lines have met recently in the continuum theory of Qu, Pan, Zhu, Jin and Roy, who build a
thermoelectric theory of piezoelectric semiconductors from a dissipation inequality and keep the Seebeck
and Peltier effects \cite{Qu2023}. That paper settles the constitutive side. What it does not do is the
plate: no free faces, no guided modes, no dispersion relation. Gupta and co-workers put a
Moore--Gibson--Thompson heat equation into a piezo-semiconductor half space with voids and studied the
size-dependent Rayleigh wave \cite{Gupta2025}; the heat operator in a semiconducting piezoelectric medium
is therefore already in print, and only the finite plate is open.

The size effect is a separate axis. Lim, Zhang and Reddy combined the nonlocal integral law with strain
gradient elasticity so that one length softens and another stiffens, and gave the equivalent differential
form together with the constitutive boundary conditions \cite{Lim2015}. Barretta and Marotti de Sciarra
explained what those boundary conditions do in a bounded body and showed that the differential law raises
the order of the problem \cite[Eqs.~(18)--(19)]{Barretta2018}. For a plate this is not a detail: the order
in the thickness coordinate goes up, and the faces need more conditions than the classical six.

Our contribution is therefore the free-face problem, the size effect and the attenuation, not a new
constitutive law. Everything in Section~\ref{sec:source} is taken from a source and cited at the equation
I read; what is new is the reduction, the counting, the parity statement and the matrix.

\section{Setting and assumptions}\label{sec:setting}

The plate occupies $-h\le z\le h$ and is unbounded in $x$ and $y$. The crystal is hexagonal, class $6mm$,
with the six-fold axis along $z$. Motion is plane strain in the $x$--$z$ plane, so $u_y=0$ and
$\partial_y=0$, and one plane wave travels along $x$,
\begin{align}
\bigl(u_x,u_z,\varphi,\Delta n,\theta\bigr)
&=\bigl(\hat u_x,\hat u_z,\hat\varphi,\hat n,\hat\theta\bigr)(z)\\
&\quad\times\exp\bigl[\ii(kx-\omega t)\bigr],
\qquad k=k_r+\ii k_i,
\label{eq:ansatz}
\end{align}
with $\omega$ real and $k_i>0$ meaning decay along $+x$. The frequency is reported through $f\,(2h)$, the
product of frequency and thickness, so that plates of different thickness sit on one axis.

The assumptions are these. Linear piezoelectricity at the reference temperature $T_0$, with the material
constants not varied with $\theta$. A small carrier perturbation, $\Delta n\ll n_0$, so that the drift
term is linearised about a uniform bias field $E_0$ along $x$ and the recombination term is linear in
$\Delta n$; the bias field is an input, not an unknown. No photoexcitation, no trapping, no surface charge
beyond what the electrical condition imposes. A small temperature rise, $|\theta|/T_0\ll1$, so the heat
equation may be written for $\theta$ with constant coefficients. A homogeneous, stress-free bulk with no
pre-stress. The two size lengths are material inputs and not fitting parameters.

\section{Notation}\label{sec:notation}

\begin{center}
\small
\begin{tabular}{p{2.1cm}p{7.6cm}p{4.3cm}}
\toprule
symbol & meaning in words & unit\\
\midrule
$2h$ & thickness of the plate; the faces are at $z=+h$ and $z=-h$ & m\\
$k=k_r+\ii k_i$ & complex wavenumber along $x$; $k_i>0$ is decay & m$^{-1}$\\
$\eta$ & exponent in the thickness direction, one value per partial wave & m$^{-1}$\\
$\omega$ & angular frequency, $2\pi f$ & rad\,s$^{-1}$\\
$u_x,u_z$ & displacements in the plane and through the thickness & m\\
$\varphi$ & electric potential; the field is $E_i=-\varphi_{,i}$ & V\\
$\Delta n$ & excess electron density over the equilibrium value $n_0$ & m$^{-3}$\\
$\theta$ & temperature rise over the reference $T_0$ & K\\
$c_{11},c_{13},c_{33},c_{44}$ & elastic moduli of class $6mm$; $c_{44}$ is the shear modulus on $xz$ & Pa\\
$e_{31},e_{33},e_{15}$ & piezoelectric constants: normal, normal along $z$, shear & C\,m$^{-2}$\\
$\epsilon_{11},\epsilon_{33}$ & dielectric permittivities in and along the plane & F\,m$^{-1}$\\
$\beta_1,\beta_3$ & thermal stress coefficients, $\beta_{ij}=c_{ijkl}\alpha_{kl}$ & Pa\,K$^{-1}$\\
$p_3$ & pyroelectric constant along the axis & C\,m$^{-2}$K$^{-1}$\\
$\rho$ & mass density & kg\,m$^{-3}$\\
$q$ & magnitude of the carrier charge, $1.602\times10^{-19}$; the electron charge is $-q$ & C\\
$b_n$ & electron mobility & m$^2$V$^{-1}$s$^{-1}$\\
$D_n$ & electron diffusion coefficient & m$^2$s$^{-1}$\\
$v_d$ & drift velocity of the bias, $v_d=b_nE_0$ & m\,s$^{-1}$\\
$S_n$ & Seebeck coefficient of the electrons & V\,K$^{-1}$\\
$\beta_{\rm th}$ & thermodiffusion coefficient, $\beta_{\rm th}=-b_nn_0S_n$ & m$^{-1}$s$^{-1}$K$^{-1}$\\
$\Pi_{\rm th}$ & Peltier coefficient, $\Pi_{\rm th}=T_0S_n$ & V\\
$\tau_R$ & recombination time; $\tau_R\to\infty$ deletes both recombination terms & s\\
$E_g$ & band gap & J\\
$\rho c_p$ & heat capacity per unit volume & J\,m$^{-3}$K$^{-1}$\\
$k_{\rm th}$ & thermal conductivity & W\,m$^{-1}$K$^{-1}$\\
$\tau_q,\tau_T$ & Moore--Gibson--Thompson phase lag, and its thermal companion in the thermoelastic theory of \cite{Quintanilla2019} & s\\
$l_{\rm nl},l_g$ & nonlocal length and gradient length; both enter as length over thickness & m\\
$\Ncal(\eta),\Gcal(\eta)$ & the two size-effect factors at the exponent $\eta$ & dimensionless\\
\bottomrule
\end{tabular}
\end{center}

Two everyday pictures, because both operators are easy to misuse. A plate is a corridor with two hard
walls: a disturbance cannot leave through a face, so it bounces and the plate guides it. The size effect is
what happens when the corridor is only a few steps wide and the walkers can feel the walls before they
reach them; a softening length lets them pass a message a little further than a local law would, a
stiffening one makes the message cost more over a short distance.

\section{Source equations}\label{sec:source}

Each equation below is written once, with the source and the equation number I read, and a line saying
whether it is taken as it stands, taken with a stated change, or ours.

Strain and field relations, as usual,
\begin{equation}
\varepsilon_{ij}=\tfrac12\bigl(u_{i,j}+u_{j,i}\bigr),\qquad E_i=-\varphi_{,i}.
\label{eq:kin}
\end{equation}
These are definitions, textbook.

The piezoelectric constitutive law without temperature is \cite[Eq.~(1)]{Zhu2023}; the same paper writes
the field-potential and strain relations used above at its Eq.~(2). I take the elastic, piezoelectric and
dielectric parts from that equation as they stand.

The thermal stress is added by Duhamel's analogy: a temperature rise acts exactly like an inelastic strain
$\alpha_{ij}\theta$, so the stress picks up $-\beta_{ij}\theta$ with
$\beta_{ij}=c_{ijkl}\alpha_{kl}$. Textbook, no equation number needed.

The pyroelectric contribution $p_i\theta$ to the electric displacement is textbook too.

The balance of momentum in the nonlocal strain gradient form, and the constitutive law that goes with it,
\begin{equation}
\bigl(1-l_{\rm nl}^2\nabla^2\bigr)t_{ij}=\sigma_{ij},\qquad
\sigma_{ij}=c_{ijkl}\bigl(1-l_g^2\nabla^2\bigr)\varepsilon_{kl}
-e_{kij}E_k-\beta_{ij}\theta ,
\label{eq:nsgt}
\end{equation}
are Eq.~(18) of Lim, Zhang and Reddy with the constitutive boundary conditions of their Eq.~(19). I did not
open that paper: the equation numbers are the ones Barretta and Marotti de Sciarra identify, in the
introduction of \cite{Barretta2018}, where they write that Lim et al.\ replaced the integral law by the
differential law Eq.~(18) and that the boundary conditions Eq.~(19) hold when the fields vanish at
infinity. Before this pair is quoted with its numbers in a manuscript, open the Lim paper and confirm (18)
on the page.

The momentum equation is used in the form in which the same operator acts on the inertia,
\begin{equation}
\bigl(1-l_{\rm nl}^2\nabla^2\bigr)\rho\,\ddot u_i=\sigma_{ij,j}.
\label{eq:mom}
\end{equation}
This form is ours, one line from \eqref{eq:nsgt}: apply the operator to the balance $t_{ij,j}=\rho\ddot
u_i$ and use $(1-l_{\rm nl}^2\nabla^2)t_{ij}=\sigma_{ij}$. The point of the form is that no derivative of
the operator ever lands on a boundary term.

Gauss's law with the mobile charge, in the linearised form used for a small carrier perturbation,
\begin{equation}
D_{i,i}=-q\,\Delta n .
\label{eq:gauss}
\end{equation}
The right-hand side is the charge added by the excess electrons, $-q\Delta n$ per unit volume. The form is
textbook electrostatics applied to a semiconductor; the same equation for a piezoelectric semiconductor
plate appears in the drift--diffusion treatment of \cite{YangZhou2005}, whose equations I read only at the
level of the abstract.

Carrier continuity with drift, diffusion, thermodiffusion and an optional recombination term,
\begin{equation}
\frac{\partial\Delta n}{\partial t}
=D_n\nabla^2\Delta n+b_nn_0\,\nabla\!\cdot\!E
+v_d\frac{\partial\Delta n}{\partial x}
+\beta_{\rm th}\nabla^2\theta-\frac{\Delta n}{\tau_R}.
\label{eq:carrier}
\end{equation}
The first three terms on the right are the standard drift--diffusion pair in the linearised form of
\cite[Eqs.~(3), (7)]{YangZhou2005}; the drift term is written with the bias velocity $v_d=b_nE_0$ rather
than solved for. The thermodiffusion term is the Soret cross-effect, textbook irreversible
thermodynamics. The recombination term is an option and is discussed below.

The heat equation of the Moore--Gibson--Thompson theory in its base form,
\begin{equation}
\Bigl(1+\tau_q\frac{\partial}{\partial t}\Bigr)
\Bigl[\rho c_p\dot\theta+T_0\beta_{ij}\dot\varepsilon_{ij}
+T_0p_i\dot E_i+q\Pi_{\rm th}\frac{\partial\Delta n}{\partial t}
+\frac{E_g\Delta n}{\tau_R}\Bigr]
=k_{\rm th}\nabla^2\theta .
\label{eq:heat}
\end{equation}
The operator in front is the Moore--Gibson--Thompson one with the conductivity rate put to zero and one
time integration applied, which is the form used in the piezo-semiconductor literature
\cite{Gupta2025}; I read the title and the abstract of that paper, not its equations, and the exact form
above is the one stated in our model document rather than copied from a page. The thermoelastic term
$T_0\beta_{ij}\dot\varepsilon_{ij}$ and the electrocaloric term $T_0p_i\dot E_i$ are textbook partners of
the thermal stress and the pyroelectric term; dropping them while keeping $\beta_{ij}$ and $p_i$ would
break reciprocity. The Peltier term is the reciprocal partner of the Soret term in \eqref{eq:carrier}, and
Qu et al.\ keep exactly this pair in a piezoelectric semiconductor continuum \cite{Qu2023}. The
recombination term $+E_g\Delta n/\tau_R$ is ours, taken from the energy balance of the electron gas; it is
an option, and it is deleted together with $-\Delta n/\tau_R$ when $\tau_R\to\infty$.

\section{Calculation A: constitutive relations in components}\label{sec:calcA}

\begin{worked}
The class is $6mm$ with the axis along $z$. The moduli are a fourth-order tensor with the axis as a
six-fold axis, so in the two-index notation $c_{11}=c_{22}$, $c_{13}=c_{23}$, $c_{33}$, $c_{44}=c_{55}$,
$c_{66}=(c_{11}-c_{12})/2$, and every element with an index $4$, $5$ or $6$ mixed with $1$, $2$ or $3$
vanishes. Plane strain kills $y$: $\varepsilon_{yy}=0$, $\partial_y=0$. What survives is
$c_{11},c_{13},c_{33},c_{44}$.
The piezoelectric tensor of $6mm$ has three independent entries, $e_{31}=e_{32}$, $e_{33}$, $e_{15}=e_{24}$.
The dielectric tensor has two, $\epsilon_{11}=\epsilon_{22}$ and $\epsilon_{33}$. The thermal stress tensor
inherits the same shape, so $\beta_{xx}=\beta_{yy}=\beta_1$ and $\beta_{zz}=\beta_3$. The pyroelectric
vector has one entry, $p_3$, along the axis.
Both components of the field survive, $E_x=-\partial_x\varphi$ and $E_z=-\partial_z\varphi$.
\end{worked}

\subsection*{The three stresses, one at a time}

\begin{worked}
Start from $\sigma_{xx}=c_{11}\varepsilon_{xx}+c_{12}\varepsilon_{yy}+c_{13}\varepsilon_{zz}
-e_{31}E_z-\beta_1\theta$. Plane strain in $y$ gives $\varepsilon_{yy}=0$, so the $c_{12}$ term drops.
The size-effect operator acts on the elastic strain energy, so it multiplies the two stiffness terms and
nothing else. Substitution: nothing left to substitute. Collect: two stiffness terms, one piezo term, one
thermal term. Result: the first line of \eqref{eq:stress}.
\end{worked}

\begin{worked}
$\sigma_{zz}=c_{13}\varepsilon_{xx}+c_{23}\varepsilon_{yy}+c_{33}\varepsilon_{zz}-e_{33}E_z
-\beta_3\theta$. With $\varepsilon_{yy}=0$ and $c_{23}=c_{13}$, substitution gives
$c_{13}\varepsilon_{xx}+c_{33}\varepsilon_{zz}$; collect the piezo and thermal terms; result: the second
line of \eqref{eq:stress}.
\end{worked}

\begin{worked}
$\sigma_{xz}=2c_{44}\varepsilon_{xz}-e_{15}E_x$. There is no $c$ term with an index $5$ mixing with $1$
or $3$, so the shear has one modulus only. Note the factor $2$: the two-index notation carries
$\sigma_{xz}=c_{55}\,2\varepsilon_{xz}$ and $c_{55}=c_{44}$. Substitution: the modulus and the piezo
coefficient. Collect: two terms. Result: the third line.
\end{worked}

\subsection*{The two displacement components, one at a time}

\begin{worked}
$D_x=e_{15}\,2\varepsilon_{xz}+\epsilon_{11}E_x$. The transverse piezo coefficient $e_{15}$ multiplies
the engineering shear $2\varepsilon_{xz}$, not $\varepsilon_{xz}$; losing this factor of two is the
commonest error in this reduction. There is no pyroelectric term in $D_x$ because the polar axis is $z$.
Result: the first of \eqref{eq:D}.
\end{worked}

\begin{worked}
$D_z=e_{31}\varepsilon_{xx}+e_{32}\varepsilon_{yy}+e_{33}\varepsilon_{zz}+\epsilon_{33}E_z+p_3\theta$.
With $\varepsilon_{yy}=0$ and $e_{32}=e_{31}$ this is $e_{31}\varepsilon_{xx}+e_{33}\varepsilon_{zz}
+\epsilon_{33}E_z+p_3\theta$. Result: the second of \eqref{eq:D}.
\end{worked}

Writing the stresses out,
\begin{align}\label{eq:stress}
\sigma_{xx}&=c_{11}\bigl(1-l_g^2\nabla^2\bigr)\varepsilon_{xx}
+c_{13}\bigl(1-l_g^2\nabla^2\bigr)\varepsilon_{zz}-e_{31}E_z-\beta_1\theta,\\
\sigma_{zz}&=c_{13}\bigl(1-l_g^2\nabla^2\bigr)\varepsilon_{xx}
+c_{33}\bigl(1-l_g^2\nabla^2\bigr)\varepsilon_{zz}-e_{33}E_z-\beta_3\theta,\\
\sigma_{xz}&=2c_{44}\bigl(1-l_g^2\nabla^2\bigr)\varepsilon_{xz}-e_{15}E_x,
\end{align}
and the electric displacement,
\begin{equation}\label{eq:D}
D_x=e_{15}\,2\varepsilon_{xz}+\epsilon_{11}E_x,\qquad
D_z=e_{31}\varepsilon_{xx}+e_{33}\varepsilon_{zz}+\epsilon_{33}E_z+p_3\theta .
\end{equation}
The two size-effect operators sit on the elastic strain energy and nowhere else. That is a modelling
choice: it keeps $D_i$ second order, so the electric problem needs no extra face condition, and it keeps
the size effect mechanical, so the extra conditions are mechanical too.

\section{Calculation B: the five governing equations in components}\label{sec:calcB}

The five equations are \eqref{eq:mom} twice, \eqref{eq:gauss}, \eqref{eq:carrier} and \eqref{eq:heat}.
Before the plane wave goes in, it is worth writing the two operators once. With
$\nabla^2=\partial_x^2+\partial_z^2$ and the plane wave \eqref{eq:ansatz},
\begin{equation}
\nabla^2\;\longrightarrow\;\eta^2-k^2\equiv L,\qquad
\Ncal(\eta)=1+l_{\rm nl}^2(k^2-\eta^2),\qquad
\Gcal(\eta)=1+l_g^2(k^2-\eta^2).
\label{eq:NG}
\end{equation}
Note the sign: $\nabla^2\to\eta^2-k^2$, so $1-l^2\nabla^2\to1+l^2(k^2-\eta^2)$, which is positive for a
gradient length when $k>\eta$ and is the stiffening factor.

\begin{worked}
Put $f=1-\ii\omega\tau_q$, the factor from the Moore--Gibson--Thompson operator, because applying
$1+\tau_q\partial_t$ to $\exp(-\ii\omega t)$ multiplies by $1-\ii\omega\tau_q$.
Put $\dot{\;\;}\to-\ii\omega$ and $\ddot{\;\;}\to-\omega^2$ everywhere.
Momentum in $x$: $\partial_x\sigma_{xx}+\partial_z\sigma_{xz}$ with $\partial_x\to\ii k$, $\partial_z\to\eta$.
Momentum in $z$: $\partial_x\sigma_{xz}+\partial_z\sigma_{zz}$.
Gauss: $\partial_xD_x+\partial_zD_z+q\Delta n=0$.
Carrier: move the whole right side of \eqref{eq:carrier} to the left.
Heat: apply \eqref{eq:NG} to the right side, which gives $k_{\rm th}L\theta$, and move it left.
\end{worked}

The five equations read $M(\eta)\,\mathbf a=\mathbf 0$ with
$\mathbf a=(\hat u_x,\hat u_z,\hat\varphi,\hat n,\hat\theta)^{T}$. The entries are given in
Section~\ref{sec:calcC} and each one is derived there.

\section{Calculation C: the plane-wave reduction}\label{sec:calcC}

First the strains and the field at the plane wave,
\begin{equation}
\varepsilon_{xx}=\ii k\hat u_x,\qquad
\varepsilon_{zz}=\eta\hat u_z,\qquad
2\varepsilon_{xz}=\eta\hat u_x+\ii k\hat u_z,\qquad
E_x=-\ii k\hat\varphi,\qquad
E_z=-\eta\hat\varphi .
\end{equation}

\begin{worked}
Substituting into the three stresses of Calculation A,
$\hat\sigma_{xx}=\Gcal(\ii k\,c_{11}\hat u_x+\eta\,c_{13}\hat u_z)+\eta\,e_{31}\hat\varphi-\beta_1\hat\theta$,
because $-e_{31}E_z=-e_{31}(-\eta\hat\varphi)=+\eta e_{31}\hat\varphi$.
Similarly $\hat\sigma_{zz}=\Gcal(\ii k\,c_{13}\hat u_x+\eta\,c_{33}\hat u_z)+\eta\,e_{33}\hat\varphi
-\beta_3\hat\theta$ and
$\hat\sigma_{xz}=\Gcal c_{44}(\eta\hat u_x+\ii k\hat u_z)+\ii k\,e_{15}\hat\varphi$, because
$-e_{15}E_x=+\ii k e_{15}\hat\varphi$.
The displacements are
$\hat D_x=e_{15}(\eta\hat u_x+\ii k\hat u_z)-\ii k\epsilon_{11}\hat\varphi$ and
$\hat D_z=\ii k e_{31}\hat u_x+\eta e_{33}\hat u_z-\eta\epsilon_{33}\hat\varphi+p_3\hat\theta$.
\end{worked}

\subsection*{Row 1: momentum in $x$}

\begin{worked}
The balance is $\Ncal\rho\ddot u_x=\partial_x\sigma_{xx}+\partial_z\sigma_{xz}$; with
$\ddot u_x=-\omega^2\hat u_x$ we write it as
$-\Ncal\rho\omega^2\hat u_x-\partial_x\hat\sigma_{xx}-\partial_z\hat\sigma_{xz}=0$, so that every term
sits on one side. Now take the three pieces in turn.
Inertia: $-\Ncal\rho\omega^2X$. This carries no $\eta$ at all, which is why the nonlocal length does not
raise the degree of the polynomial.
First divergence: $-\partial_x\hat\sigma_{xx}=-\ii k\bigl[\Gcal(\ii kc_{11}X+\eta c_{13}Z)
+\eta e_{31}\Phi-\beta_1\Theta\bigr]
=\Gcal c_{11}k^2X-\ii k\eta\Gcal c_{13}Z-\ii k\eta e_{31}\Phi+\ii k\beta_1\Theta$.
Second divergence: $-\partial_z\hat\sigma_{xz}=-\eta\bigl[\Gcal c_{44}(\eta X+\ii kZ)
+\ii k e_{15}\Phi\bigr]=-\Gcal c_{44}\eta^2X-\ii k\eta\Gcal c_{44}Z-\ii k\eta e_{15}\Phi$.
Collect the coefficient of $X$: $-\Ncal\rho\omega^2+\Gcal c_{11}k^2-\Gcal c_{44}\eta^2
=-\Ncal\rho\omega^2+\Gcal(c_{11}k^2-c_{44}\eta^2)$. That is $M_{11}$.
Collect $Z$: $-\ii k\eta\Gcal c_{13}-\ii k\eta\Gcal c_{44}=-\ii k\eta\Gcal(c_{13}+c_{44})$.
That is $M_{12}$.
Collect $\Phi$: $-\ii k\eta e_{31}-\ii k\eta e_{15}=-\ii k\eta(e_{31}+e_{15})$. That is $M_{13}$.
Collect $\Theta$: $+\ii k\beta_1$. That is $M_{15}$.
The carrier amplitude $\hat n$ appears nowhere in the momentum balance, so $M_{14}=0$.
Result: row 1 of \eqref{eq:M1}.
\end{worked}

\begin{equation}\label{eq:M1}
M_{11}=-\Ncal\rho\omega^2+\Gcal\bigl(c_{11}k^2-c_{44}\eta^2\bigr),\quad
M_{12}=-\ii k\eta\,\Gcal(c_{13}+c_{44}),\quad
M_{13}=-\ii k\eta\,(e_{31}+e_{15}),\quad
M_{15}=\ii k\,\beta_1 .
\label{eq:M1}
\end{equation}

\subsection*{Row 2: momentum in $z$}

\begin{worked}
The balance is $\Ncal\rho\ddot u_z=\partial_x\sigma_{xz}+\partial_z\sigma_{zz}$, written as
$-\Ncal\rho\omega^2Z-\partial_x\hat\sigma_{xz}-\partial_z\hat\sigma_{zz}=0$.
Inertia: $-\Ncal\rho\omega^2Z$.
First divergence: $-\partial_x\hat\sigma_{xz}=-\ii k\bigl[\Gcal c_{44}(\eta X+\ii kZ)
+\ii k e_{15}\Phi\bigr]=-\ii k\eta\Gcal c_{44}X+\Gcal c_{44}k^2Z+e_{15}k^2\Phi$, because
$-\ii k\cdot\ii k=+k^2$.
Second divergence: $-\partial_z\hat\sigma_{zz}=-\eta\bigl[\Gcal(\ii kc_{13}X+\eta c_{33}Z)
+\eta e_{33}\Phi-\beta_3\Theta\bigr]
=-\ii k\eta\Gcal c_{13}X-\Gcal c_{33}\eta^2Z-e_{33}\eta^2\Phi+\eta\beta_3\Theta$.
Collect $X$: $-\ii k\eta\Gcal(c_{44}+c_{13})$. That is $M_{21}$, and it equals $M_{12}$, as it must:
the two momentum rows come from one symmetric operator.
Collect $Z$: $-\Ncal\rho\omega^2+\Gcal c_{44}k^2-\Gcal c_{33}\eta^2
=-\Ncal\rho\omega^2+\Gcal(c_{44}k^2-c_{33}\eta^2)$. That is $M_{22}$.
Collect $\Phi$: $e_{15}k^2-e_{33}\eta^2$. That is $M_{23}$.
Collect $\Theta$: $+\eta\beta_3$. That is $M_{25}$, and it is odd in $\eta$, unlike $M_{15}$, which is
even; the difference is that $\beta_1$ enters through $\partial_x$ and $\beta_3$ through $\partial_z$.
$M_{24}=0$. Result: row 2 of \eqref{eq:M2}.
\end{worked}

\begin{equation}\label{eq:M2}
M_{21}=-\ii k\eta\,\Gcal(c_{13}+c_{44}),\quad
M_{22}=-\Ncal\rho\omega^2+\Gcal\bigl(c_{44}k^2-c_{33}\eta^2\bigr),\quad
M_{23}=e_{15}k^2-e_{33}\eta^2,\quad
M_{25}=\eta\,\beta_3 .
\label{eq:M2}
\end{equation}

\subsection*{Row 3: Gauss law}

\begin{worked}
Gauss' law with the free charge kept is $\partial_xD_x+\partial_zD_z+q\Delta n=0$; there is
no external charge, and the carrier density is the only free charge in the problem.
First term: $\partial_x\hat D_x=\ii k\bigl[e_{15}(\eta X+\ii kZ)-\ii k\epsilon_{11}\Phi\bigr]
=\ii k\eta e_{15}X-e_{15}k^2Z+\epsilon_{11}k^2\Phi$.
Second term: $\partial_z\hat D_z=\eta\bigl[\ii ke_{31}X+\eta e_{33}Z-\eta\epsilon_{33}\Phi
+p_3\Theta\bigr]=\ii k\eta e_{31}X+e_{33}\eta^2Z-\epsilon_{33}\eta^2\Phi+\eta p_3\Theta$.
Free charge: $+q\hat n$.
Collect $X$: $\ii k\eta(e_{15}+e_{31})$. That is $M_{31}$, and it is $-M_{13}$, the sign that makes the
piezoelectric part of the operator skew.
Collect $Z$: $e_{33}\eta^2-e_{15}k^2$. That is $M_{32}=-M_{23}$, for the same reason.
Collect $\Phi$: $\epsilon_{11}k^2-\epsilon_{33}\eta^2$. That is $M_{33}$.
Collect $\hat n$: $q$. That is $M_{34}$, the one entry of the matrix with no $\eta$ and no $k$.
Collect $\Theta$: $\eta p_3$. That is $M_{35}$, the pyroelectric entry.
Result: row 3 of \eqref{eq:M3}.
\end{worked}

\begin{equation}\label{eq:M3}
M_{31}=\ii k\eta(e_{15}+e_{31}),\quad
M_{32}=e_{33}\eta^2-e_{15}k^2,\quad
M_{33}=\epsilon_{11}k^2-\epsilon_{33}\eta^2,\quad
M_{34}=q,\quad
M_{35}=\eta\,p_3 .
\label{eq:M3}
\end{equation}

\subsection*{Row 4: carrier continuity}

\begin{worked}
The carrier balance carries diffusion, drift, thermodiffusion and recombination:
$\partial_t\Delta n=D_n\nabla^2\Delta n-b_nn_0\nabla^2\varphi+v_d\,\partial_x\Delta n
-\Delta n/\tau_R+\beta_{\rm th}\nabla^2\theta$. The sign of the drift term is the sign of the bias
velocity; we take the bias along $+x$ and write the term on the right, which puts $-\ii kv_d$ in the
matrix entry. Now substitute $\partial_t\to-\ii\omega$, $\nabla^2\to L$, $\partial_x\to\ii k$:
$-\ii\omega\hat n=D_nL\hat n-b_nn_0L\Phi+\ii kv_d\hat n-\hat n/\tau_R+\beta_{\rm th}L\Theta$.
Move everything to the left:
$\bigl(-\ii\omega-D_nL-\ii kv_d+1/\tau_R\bigr)\hat n+b_nn_0L\,\Phi-\beta_{\rm th}L\,\Theta=0$.
Collect $\Phi$: $b_nn_0L=b_nn_0(\eta^2-k^2)$. That is $M_{43}$.
Collect $\hat n$: $-\ii\omega-D_nL-\ii kv_d+1/\tau_R$. That is $M_{44}$; the three real pieces are
the diffusion rate $D_n(k^2-\eta^2)$, the recombination rate $1/\tau_R$ and the drift $-\ii kv_d$,
which is the only term that makes $M_{44}$ complex at $\omega=0$.
Collect $\Theta$: $-\beta_{\rm th}L$. That is $M_{45}$, even in $\eta$.
$M_{41}=M_{42}=0$: the carrier equation has no mechanical driving term. Result: row 4 of \eqref{eq:M4}.
\end{worked}

\begin{equation}\label{eq:M4}
M_{43}=b_nn_0\,L,\quad
M_{44}=-\ii\omega-D_nL-\ii k\,v_d+\frac{1}{\tau_R},\quad
M_{45}=-\beta_{\rm th}L .
\label{eq:M4}
\end{equation}

\subsection*{Row 5: the heat equation}

\begin{worked}
The heat equation is of Moore--Gibson--Thompson type, which means the flux law is
$q_i+\tau_q\dot q_i=-k_{\rm th}\theta_{,i}$ and the balance is $\rho c_p\dot\theta+q_{i,i}=S$
with the sources $S=T_0\beta_1\dot\varepsilon_{xx}+T_0\beta_3\dot\varepsilon_{zz}
+T_0p_3\dot E_z+q\Pi_{\rm th}\dot{\Delta n}+(E_g/\tau_R)\Delta n$.
Apply $1+\tau_q\partial_t$ to the balance and use the flux law for the divergence term:
$\rho c_p(\dot\theta+\tau_q\ddot\theta)+\bigl(1+\tau_q\partial_t\bigr)q_{i,i}
=\bigl(1+\tau_q\partial_t\bigr)S$ and $\bigl(1+\tau_q\partial_t\bigr)q_{i,i}=-k_{\rm th}\nabla^2\theta$,
so that the conduction term carries no $f$ while every source does. With
$f=1-\ii\omega\tau_q$ and the sources on the left,
$f\bigl[T_0\beta_1\dot\varepsilon_{xx}+T_0\beta_3\dot\varepsilon_{zz}+T_0p_3\dot E_z
+q\Pi_{\rm th}\dot{\Delta n}+(E_g/\tau_R)\Delta n\bigr]
+k_{\rm th}L\,\Theta+\ii\omega f\rho c_p\,\Theta=0$.
Now the five substitutions, one per line.
$\dot\varepsilon_{xx}=-\ii\omega\,\ii kX=\omega kX$, because $(-\ii)(\ii)=+1$. With the prefactor
$fT_0\beta_1$ this gives $M_{51}=fT_0\omega k\beta_1$, even in $\eta$.
$\dot\varepsilon_{zz}=-\ii\omega\,\eta Z$, so the term is $-\ii\omega fT_0\beta_3\eta Z$ and
$M_{52}=-\ii\omega fT_0\beta_3\eta$.
$E_z=-\eta\Phi$, so $\dot E_z=+\ii\omega\eta\Phi$ and $M_{53}=\ii\omega fT_0p_3\eta$. Both
$M_{52}$ and $M_{53}$ are odd, and both fit the parity matching of Section~\ref{sec:calcD}.
$\dot{\Delta n}=-\ii\omega\hat n$ and $\Delta n$ enters undifferentiated, so the carrier column
reads $f(-\ii\omega q\Pi_{\rm th}+E_g/\tau_R)$. That is $M_{54}$, even in $\eta$, and it is one
of the two entries that break the split.
$\Theta$ collects $-k_{\rm th}L$ from the conduction and $-\ii\omega f\rho c_p$ from the balance, so
$M_{55}=-\ii\omega f\rho c_p-k_{\rm th}L$. Setting $\tau_q=0$ makes $f=1$ and returns the Fourier
law; that is the limit used in Calculation~H.
Result: row 5 of \eqref{eq:M5}.
\end{worked}

\begin{align}
\label{eq:M5} \\
M_{51}&==f\,T_0\omega k\,\beta_1, \\
M_{52}&==-\ii\omega f\,T_0\beta_3\eta, \\
M_{53}&==\ii\omega f\,T_0p_3\eta, \\
M_{54}&==f\Bigl(-\ii\omega q\Pi_{\rm th}+\frac{E_g}{\tau_R}\Bigr), \\
M_{55}&==-\ii\omega f\rho c_p-k_{\rm th}L .
\end{align}

All other entries are zero. The matrix is
\begin{equation}
M(\eta)=
\begin{pmatrix}
M_{11}&M_{12}&M_{13}&0&M_{15}\\
M_{21}&M_{22}&M_{23}&0&M_{25}\\
M_{31}&M_{32}&M_{33}&M_{34}&M_{35}\\
0&0&M_{43}&M_{44}&M_{45}\\
M_{51}&M_{52}&M_{53}&M_{54}&M_{55}
\end{pmatrix}.
\label{eq:Mmat}
\end{equation}
The pattern of zeros is worth reading: the carrier talks to the mechanical and electric fields only
through Gauss's law and to the temperature only through the Soret and Peltier terms, and it has no
inertia. That is why the carrier partial waves are short and diffusive, and why the matrix is not
symmetric.

\section{Calculation D: the characteristic equation}\label{sec:calcD}

Non-trivial amplitudes need $\det M(\eta)=0$. The determinant is a polynomial in $\eta$.

\begin{worked}
Count the degree row by row, taking the highest power of $\eta$ in each entry. In row 1 the entry
$M_{11}$ contains $\Gcal$, which is quadratic in $\eta$, multiplied by $c_{44}\eta^2$, so it is of degree
four; $M_{12}$ and $M_{13}$ are degree three; $M_{15}$ is degree one. In row 2 the same count gives four,
four (in $M_{22}$), two and one. Row 3 has degree one, two, two, zero and one. Row 4 has zero, zero, two,
two and two. Row 5 has zero ($M_{51}$ carries no $\eta$), one, one, zero and two.
The determinant is a sum of products, one entry from each row and each column, so its degree is the
largest sum over the permutations. Taking the diagonal gives $4+4+2+2+2=14$. Any permutation that gives
up $M_{11}$ or $M_{22}$ loses at least two and gains at most two elsewhere, so 14 is the maximum, and the
coefficient of $\eta^{14}$ does not vanish because it multiplies $l_g^4c_{44}c_{33}$ and the diagonal
entry $M_{44}M_{55}$.
\end{worked}

With $l_g\ne0$ the polynomial has degree $14$. With $l_g=0$ the factors $\Gcal$ become unity, $M_{11}$ and
$M_{22}$ drop to degree two, and the same counting gives $2+2+2+2+2=10$. The nonlocal length alone does not
raise the order, because $\Ncal$ sits on the inertia, which carries no $\eta$. Without any size effect the
degree is 10 as well.

Fourteen roots means fourteen partial waves, and the free-face problem needs fourteen conditions, which is
what Section~\ref{sec:calcF} supplies.

\subsection*{Parity}

The parity of $\det M$ decides whether the plate modes split into a symmetric and an antisymmetric family,
and the answer here is not the classical one.

Write the parity of every entry first. Each entry is a monomial in $\eta$ times a constant, so its parity
is decided by the power of $\eta$ and by nothing else:
\begin{equation}
\begin{array}{c|ccccc}
 & \hat u_x & \hat u_z & \hat\varphi & \hat n & \hat\theta\\\hline
\text{row 1} & + & - & - & \cdot & +\\
\text{row 2} & - & + & + & \cdot & -\\
\text{row 3} & - & + & + & + & -\\
\text{row 4} & \cdot & \cdot & + & + & +\\
\text{row 5} & + & - & - & + & +
\end{array}
\label{eq:paritytable}
\end{equation}
with $+$ for even, $-$ for odd and $\cdot$ for a zero entry.

\begin{worked}
Now look for signs. The rule is $M_{ij}(-\eta)=\sigma_i\tau_jM_{ij}(\eta)$ for every non-zero entry, the
row signs $\sigma_i$ and the column signs $\tau_j$ being $\pm1$.
Set $\sigma_1=+1$, which is free because only the products matter. From $M_{11}$ even, $\tau_1=+1$; from
$M_{12}$ odd and $M_{13}$ odd, $\tau_2=\tau_3=-1$; from $M_{15}$ even, $\tau_5=+1$.
Row 2: $M_{21}$ odd gives $\sigma_2=-\tau_1^{-1}$, so $\sigma_2=-1$. Then $M_{22}$ even needs
$\sigma_2\tau_2=+1$, which holds; $M_{23}$ even needs $\sigma_2\tau_3=+1$, which holds; $M_{25}$ odd needs
$\sigma_2\tau_5=-1$, which holds.
Row 3: $M_{31}$ odd gives $\sigma_3=-1$. Then $M_{32}$ even, $M_{33}$ even and $M_{35}$ odd all hold
against $\tau_2=\tau_3=-1$, $\tau_5=+1$. The entry $M_{34}=q$ is even and gives $\sigma_3\tau_4=+1$, so
$\tau_4=-1$.
Row 4: $M_{43}$ even gives $\sigma_4=-\tau_3^{-1}$, so $\sigma_4=-1$; $M_{44}$ even gives
$\sigma_4\tau_4=+1$, which holds. The entry $M_{45}=-\beta_{\rm th}L$ is even and gives
$\sigma_4\tau_5=+1$, but $\sigma_4=-1$ and $\tau_5=+1$. this one fails.
Row 5: $M_{51}$ even gives $\sigma_5=+1$; $M_{52}$ odd and $M_{53}$ odd hold against $\tau_2=\tau_3=-1$;
$M_{55}$ even holds. The entry $M_{54}=f(-\ii\omega q\Pi_{\rm th}+E_g/\tau_R)$ is even and gives
$\sigma_5\tau_4=+1$, but $\tau_4=-1$. this one fails.
The assignment is forced up to one overall sign, so these two are the only failures.
\end{worked}

The conclusion is short. Two entries refuse the matching, $M_{45}=-\beta_{\rm th}L$ and
$M_{54}=f(-\ii\omega q\Pi_{\rm th}+E_g/\tau_R)$, and both of them are even in $\eta$ where the matching
needs them odd. So $\det M$ is even, and the plate modes split into a symmetric and an antisymmetric
family, if and only if
\begin{equation}
\beta_{\rm th}=0,\qquad \Pi_{\rm th}=0,\qquad E_g=0
\end{equation}
together. If any one of the three is present the polynomial carries odd powers, the fourteen roots do not
pair, and the determinant has to be taken over all fourteen at once.

The pyroelectric constant has nothing to do with it. Its two entries, $M_{35}=\eta p_3$ in the Gauss row
and $M_{53}=\ii\omega fT_0p_3\eta$ in the heat row, are both odd, and both fit the matching; setting
$p_3=0$ does not restore the pairing and leaving it in does not break it.

The physical reason for the failure is worth one sentence. Under $z\to-z$ the amplitudes of one family
transform by the column signs just found: $\hat u_x$ and $\hat\theta$ are even in $z$, while $\hat u_z$,
$\hat\varphi$ and $\hat n$ are odd. Thermodiffusion drives the carrier density from the temperature
gradient, and an even $\theta$ drives an even $\Delta n$ through the even operator $L$; the family needs
$\Delta n$ odd, so the coupling is inconsistent with the split. The Peltier and recombination terms fail
the same way in the other direction, feeding an odd $\Delta n$ into the equation for an even $\theta$.
Remove those three terms and the carrier and the temperature stop talking to each other through even
operators, and the split comes back.

A practical consequence follows. When the split holds, two smaller matrices may be built, one of order
seven for each family, and each may be factored separately. When it fails, the full fourteen-by-fourteen
matrix is the only object available. The counting of Section~\ref{sec:calcF} is unaffected either way:
fourteen conditions in, fourteen amplitudes out.

\section{Calculation E: the solution through the thickness}\label{sec:calcE}

Let $\eta_1,\dots,\eta_{14}$ be the roots of the characteristic equation. For each root the amplitude
vector is a null vector of $M(\eta_m)$.

\begin{worked}
For a matrix that is singular at $\eta_m$, the adjugate is not zero and every non-zero column of it is a
null vector. Write $C(\eta)$ for the cofactor matrix, $C_{ij}=(-1)^{i+j}\mathcal{M}_{ji}$ where
$\mathcal M_{ji}$ is the minor obtained by deleting row $j$ and column $i$. Then
$\sum_j M_{ij}(\eta_m)C_{jk}(\eta_m)=\delta_{ik}\det M(\eta_m)=0$, so column $k$ of $C$ is annihilated
by $M(\eta_m)$. Choose one column index $\alpha$ that is not identically zero and set
$a_{m,j}=C_{\alpha j}(\eta_m)$.
In practice the cleanest choice uses the minor of the $4\times4$ block obtained by deleting one row and
one column, evaluated at the root; the resulting vector is defined up to one scalar, which is absorbed into
the unknown amplitude $A_m$.
\end{worked}

The general solution is
\begin{equation}
\begin{pmatrix}u_x\\u_z\\\varphi\\\Delta n\\\theta\end{pmatrix}
=\sum_{m=1}^{14}A_m
\begin{pmatrix}X_m\\Z_m\\\Phi_m\\\nu_m\\\Theta_m\end{pmatrix}
e^{\eta_m z}\,e^{\ii(kx-\omega t)},
\label{eq:solution}
\end{equation}
with fourteen constants $A_m$ fixed by the fourteen face conditions. When $l_g=0$ the sum has ten terms.
This partial-wave construction is the standard one for a layer; the reduction to $M(\eta)$ is ours.

\section{Calculation F: the seven face conditions}\label{sec:calcF}

At each face seven conditions hold, so fourteen in all, one per unknown $A_m$. All are homogeneous. For a
root $\eta_m$ write $\Ncal_m=\Ncal(\eta_m)$ and $\Gcal_m=\Gcal(\eta_m)$.

\emph{The two tractions.} On a free face of a strain-gradient continuum the traction is not just
$\sigma_{ij}n_j$: the surface-gradient terms contribute, and the standard statement is
\cite[Sec.~2, Eq.~(14)]{Charal2019}, whose Sec.~2 condition (b) also names the double traction
$R=n\cdot\mu\cdot n$. The older statement is Mindlin's \cite{Mindlin1964}, which I did not read. In the
nonlocal part the operator is inverted root by root, because $\Ncal^{-1}$ acts on $e^{\eta z}$ as
$1/\Ncal(\eta_m)$.

\begin{worked}
The normal traction on the face $z=+h$ has the physical part $\sigma_{zz}$ and the surface-gradient part,
which for our two-dimensional setting is the tangential derivative of the couple-like quantity. Writing
the strain-gradient contribution with the gradient length, the two face rows become
$\Ncal^{-1}[\hat\sigma_{zz}-l_g^2\ii k\eta\,c_{44}(\eta X+\ii kZ)]$ and
$\Ncal^{-1}[\hat\sigma_{xz}-l_g^2\ii k\eta(\ii k c_{11}X+\eta c_{13}Z)]$. The two double tractions are
the quantities the gradient term does work on, proportional to $l_g^2$ and vanishing when $l_g=0$.
\end{worked}

The seven row functions, with $\eta=\eta_m$ and the amplitudes $(X_m,Z_m,\Phi_m,\nu_m,\Theta_m)$,
\begin{align}
b_1&==\frac{1}{\Ncal_m}\Bigl[\hat\sigma_{zz}-l_g^2\ii k\eta\,c_{44}(\eta X_m \\
\quad +\ii kZ_m)\Bigr], \\
b_2&==\frac{1}{\Ncal_m}\Bigl[\hat\sigma_{xz}-l_g^2\ii k\eta\,(\ii k c_{11}X_m \\
\quad +\eta c_{13}Z_m)\Bigr], \\
b_3&==\frac{l_g^2}{\Ncal_m}\,\eta\,c_{44}(\eta X_m+\ii kZ_m), \\
b_4&==\frac{l_g^2}{\Ncal_m}\,\eta\,(\ii k c_{13}X_m+\eta c_{33}Z_m), \\
b_5&==\ii k e_{31}X_m+\eta e_{33}Z_m-\eta\epsilon_{33}\Phi_m+p_3\Theta_m, \\
b_6&==-D_n\eta\,\nu_m+b_nn_0\eta\,\Phi_m-\beta_{\rm th}\eta\,\Theta_m, \\
b_7&==\eta\,\Theta_m .
\end{align}
The first two are the normal and shear tractions with the surface-gradient term and the nonlocal operator
inverted; the third and fourth are the double tractions; the fifth is the open-circuit condition; the
sixth is zero carrier flux; the seventh is an adiabatic face.

\subsection*{The seven rows, one at a time}

\begin{worked}
Row $b_1$. On a face of a strain-gradient solid the normal traction is not $\sigma_{zz}$ by itself; the
surface-gradient term is the tangential divergence of the double traction, which brings one more
$\partial_x$ and therefore one more $\ii k$ and one more $\eta$. Substitution: insert
$\hat\sigma_{zz}$ from Calculation~C and the amplitudes $X_m,Z_m,\Phi_m,\Theta_m$; the gradient term
contributes $-l_g^2\ii k\eta\,c_{44}(\eta X_m+\ii kZ_m)$. Collect: four amplitudes, and $\hat n$ does not
appear. Result: the first row function, with $\Ncal_m^{-1}$ outside because the nonlocal operator is
inverted root by root.
\end{worked}

\begin{worked}
Row $b_2$. The same statement for the shear traction: the physical part is $\hat\sigma_{xz}$ and the
surface-gradient part is $-l_g^2\ii k\eta(\ii kc_{11}X_m+\eta c_{13}Z_m)$, with $c_{11}$ and $c_{13}$
where $c_{44}$ and $c_{33}$ stood in $b_1$, because the shear traction is built from the gradient of the
normal strains rather than of the shear strain. Result: the second row function.
\end{worked}

\begin{worked}
Rows $b_3$ and $b_4$ are the double tractions, the quantities the gradient term does work against. They
are $l_g^2$ times the normal derivative of the stress that the corresponding strain would produce:
$l_g^2\eta\,c_{44}(\eta X_m+\ii kZ_m)$ for the shear pair and $l_g^2\eta(\ii kc_{13}X_m+\eta c_{33}Z_m)$
for the normal pair. Both vanish with $l_g$, which is why the five-row set is enough when $l_g=0$.
Result: the third and fourth row functions.
\end{worked}

\begin{worked}
Row $b_5$ is $\hat D_z$ at the face, with the amplitudes put in. Substitution:
$\hat D_z=\ii ke_{31}X_m+\eta e_{33}Z_m-\eta\epsilon_{33}\Phi_m+p_3\Theta_m$. Collect: four terms. Setting
it to zero is the open-circuit condition, because the exterior carries no free charge and no normal
displacement may leave the face. The alternative, a metallised face held at a fixed potential, replaces
the row by $\Phi_m=0$; the two give different speeds and the difference between them is the
electromechanical coupling, which is the subject of Calculation~H.
\end{worked}

\begin{worked}
Row $b_6$ is the particle flux
$\Gamma_z=-D_n\partial_z\Delta n-b_nn_0E_z-\beta_{\rm th}\partial_z\theta$, set to zero: no carrier
crosses the face. Substitution: $\partial_z\to\eta$, $E_z=-\eta\Phi_m$, $\Delta n\to\nu_m$,
$\theta\to\Theta_m$. Collect: $-D_n\eta\nu_m+b_nn_0\eta\Phi_m-\beta_{\rm th}\eta\Theta_m$. That is the
drift--diffusion current of \cite[Eqs.~(3), (7)]{YangZhou2005} with the Soret term added. Result: the
sixth row function. Note that every term carries one $\eta$, so the row vanishes identically at
$\eta=0$; that root has to be handled separately if it arises.
\end{worked}

\begin{worked}
Row $b_7$ is $\partial_z\theta=0$, an adiabatic face. Substitution: $\eta\Theta_m$. Result: the seventh
row function. Because the carrier flux vanishes on the same face, the Peltier part of the heat current
vanishes there too and no cross term survives. An isothermal face replaces the row by $\Theta_m$; a
face with a heat transfer coefficient replaces it by $\eta\Theta_m+\lambda\Theta_m$ with $\lambda$ the
surface conductance over the conductivity. Result: the seventh row function.
\end{worked}

\section{Calculation G: the dispersion matrix}\label{sec:calcG}

Assemble one row per condition and face and one column per root,
\begin{equation}
\mathbf B(k,\omega)=
\begin{pmatrix}
b_r(\eta_m)\,e^{\eta_mh}\\
b_r(\eta_m)\,e^{-\eta_mh}
\end{pmatrix}_{r=1\dots7,\;m=1\dots14},
\qquad \mathbf B\,\mathbf A=\mathbf 0,
\qquad \mathbf A=(A_1,\dots,A_{14})^{T},
\label{eq:Bmat}
\end{equation}
the upper seven rows at $z=+h$ and the lower seven at $z=-h$. The dispersion relation is
$\det\mathbf B(k,\omega)=0$, solved for complex $k$ at real $\omega$; the attenuation per unit length is
$k_i$. When $l_g=0$ use the five rows $b_1,b_2,b_5,b_6,b_7$ with ten roots and a $10\times10$ matrix.

One remark on conditioning, which is not an algorithm but a fact about the matrix: columns belonging to
roots with large $|\operatorname{Re}\eta_m|$, and the thermal and carrier boundary layers are of that kind,
grow like $e^{|\operatorname{Re}\eta_m|h}$ in the upper block and shrink in the lower. Scaling column $m$
by $e^{-|\operatorname{Re}\eta_m|h}$ multiplies the determinant by a non-zero number and leaves its zeros
where they are.

\section{Calculation H: two hand limits}\label{sec:calcH}

\subsection*{The elastic limit and the Rayleigh--Lamb relation}

Switch off every coupling: $e=0$, $p_3=0$, $\beta=0$, $k_{\rm th}=0$ with the heat row dropped, and the
carrier terms dropped. That is an elastic plate.

\begin{worked}
What remains is the two momentum equations with $M_{11}=-\rho\omega^2+c_{11}k^2-c_{44}\eta^2$,
$M_{22}=-\rho\omega^2+c_{44}k^2-c_{33}\eta^2$ and $M_{12}=M_{21}=-\ii k\eta(c_{13}+c_{44})$, the
size-effect factors being unity when both lengths vanish. Put $A=\rho\omega^2$ and $x=\eta^2$ and
substitute into the determinant of that block:
$\bigl(c_{11}k^2-c_{44}x-A\bigr)\bigl(c_{44}k^2-c_{33}x-A\bigr)-\bigl(-\ii k\eta(c_{13}+c_{44})\bigr)^2=0$.
The last term is $+k^2x(c_{13}+c_{44})^2$, because $(-\ii)^2=-1$ and there is a minus in front. Collect
in powers of $x$:
$c_{44}c_{33}x^2-x\Bigl[c_{11}(c_{11}k^2-A)+c_{44}(c_{44}k^2-A)-k^2(c_{13}+c_{44})^2\Bigr]
+(c_{11}k^2-A)(c_{44}k^2-A)=0$, where in the constant term I have used $c_{33}=c_{11}$, which is the
isotropic identity in this plane.
Now ask when the quadratic factors as $c_{11}c_{44}\bigl(x-(k^2-A/c_{11})\bigr)\bigl(x-(k^2-A/c_{44})\bigr)$.
Compare the coefficient of $x$: the factorised form gives $2c_{11}c_{44}k^2-A(c_{11}+c_{44})$, the
expanded form gives $k^2\bigl(c_{11}^2+c_{44}^2-(c_{13}+c_{44})^2\bigr)-A(c_{11}+c_{44})$. They agree if
and only if $c_{11}^2+c_{44}^2-(c_{13}+c_{44})^2=2c_{11}c_{44}$, that is
$(c_{11}-c_{44})^2=(c_{13}+c_{44})^2$, that is $c_{11}=c_{13}+2c_{44}$.
Substitution of the isotropic values $c_{11}=\lambda+2\mu$, $c_{13}=\lambda$, $c_{44}=\mu$ gives
$\lambda+2\mu=\lambda+2\mu$. So the factorisation is an identity for an isotropic solid and for no other;
a transversely isotropic plate does not factor this way and the two branches are not a clean
longitudinal and transverse pair.
Result: $x=k^2-\omega^2/c_L^2$ and $x=k^2-\omega^2/c_T^2$ with $\rho c_L^2=c_{11}$ and
$\rho c_T^2=c_{44}$ in this two-dimensional reduction, which is the usual longitudinal and transverse
pair. The face rows $b_1,b_2$ reduce to $\hat\sigma_{zz}=0$ and $\hat\sigma_{xz}=0$, and the partial waves
come in pairs $\pm\eta$.
\end{worked}

\begin{worked}
The four conditions then collapse to the two classical Rayleigh--Lamb equations. Take the amplitudes of
the longitudinal pair as $\pm1$ in $X$ with $Z$ following from row 1, and those of the transverse pair as
$\pm1$ in $Z$ with $X$ following; write the four rows, add and subtract the two face blocks, and the
system splits into a symmetric part and an antisymmetric part. Substitution of $\eta=\pm\ii p$ and
$\eta=\pm\ii q$ with $p^2=\omega^2/c_L^2-k^2$ and $q^2=\omega^2/c_T^2-k^2$ turns the exponentials into
sines and cosines. Collect and result:
$\bigl(q^2-k^2\bigr)^2\cos(ph)\sin(qh)/q+4k^2p\sin(ph)\cos(qh)=0$ for the symmetric family and the same
expression with sine and cosine exchanged for the antisymmetric one. I write the result rather than every
line of the reduction; the scholar should carry out that reduction and check that the two equations are
the classical ones.
\end{worked}

The result to check is that with all couplings off and $l_g=l_{\rm nl}=0$ the dispersion of the plate is
the classical Rayleigh--Lamb one, both families, and that the lowest antisymmetric branch tends to the
Rayleigh surface-wave speed of the half space as $f(2h)$ grows.

\subsection*{The thin-plate open-circuit speed}

With the piezoelectric coupling on and the plate thin, the lowest symmetric mode is a plate extension.

\begin{worked}
Take $k\to0$ and keep only the in-plane motion. The lowest symmetric branch is then a plate extension:
$\varepsilon_{xx}$ is the only strain that survives to leading order, and the plate is thin enough that
the two face conditions hold through the thickness, not only at the faces. So impose
$\sigma_{zz}=0$ and $D_z=0$ inside the plate, and eliminate $\varepsilon_{zz}$ and $E_z$ in favour of
$\varepsilon_{xx}$.
Substitution: the two conditions form a $2\times2$ system,
$\begin{pmatrix}c_{33}&-e_{33}\\ e_{33}&\epsilon_{33}\end{pmatrix}
\begin{pmatrix}\varepsilon_{zz}\\ E_z\end{pmatrix}
=-\varepsilon_{xx}\begin{pmatrix}c_{13}\\ e_{31}\end{pmatrix}$.
Its determinant is $\Delta_c=c_{33}\epsilon_{33}+e_{33}^2$.
Collect: solving,
$\varepsilon_{zz}=-\varepsilon_{xx}(c_{13}\epsilon_{33}+e_{33}e_{31})/\Delta_c$ and
$E_z=-\varepsilon_{xx}(c_{33}e_{31}-e_{33}c_{13})/\Delta_c$.
Substitute both into $\sigma_{xx}=c_{11}\varepsilon_{xx}+c_{13}\varepsilon_{zz}-e_{31}E_z$:
$\sigma_{xx}=\varepsilon_{xx}\Bigl[c_{11}-\dfrac{c_{13}^2\epsilon_{33}+2c_{13}e_{33}e_{31}-c_{33}e_{31}^2}
{c_{33}\epsilon_{33}+e_{33}^2}\Bigr]$.
Result: the effective in-plane modulus is
$\bar c_{11}=c_{11}-\dfrac{c_{13}^2\epsilon_{33}+2c_{13}e_{33}e_{31}-c_{33}e_{31}^2}
{c_{33}\epsilon_{33}+e_{33}^2}$,
or, in the equivalent form that is easier to check,
$\bar c_{11}=c_{11}-\dfrac{c_{13}^2}{c_{33}}
+\dfrac{\bigl(e_{31}-e_{33}c_{13}/c_{33}\bigr)^2}{\epsilon_{33}+e_{33}^2/c_{33}}$,
the two forms being equal after multiplication by $c_{33}$. The first correction is the relief of the
thickness constraint, the second is the piezoelectric stiffening on top of it.
The extensional speed follows from $\rho\omega^2=\bar c_{11}k^2$, so
$v_{\rm open}=\sqrt{\bar c_{11}/\rho}$.
For the shorted face, $E_z=0$ instead of $D_z=0$; then $\varepsilon_{zz}=-(c_{13}/c_{33})\varepsilon_{xx}$
alone and $\bar c_{11}^{\rm short}=c_{11}-c_{13}^2/c_{33}$, giving
$v_{\rm short}=\sqrt{\bar c_{11}^{\rm short}/\rho}$. The relative difference
$(v_{\rm open}^2-v_{\rm short}^2)/v_{\rm short}^2$ is the squared electromechanical coupling of the plate,
which is the quantity a measurement of the two speeds returns.
Two remarks, both worth the scholar's attention. First, the relief term $c_{13}^2/c_{33}$ is not optional:
a plate is free in $z$ and the extensional modulus of a free plate is not $c_{11}$. Writing
$\bar c_{11}=c_{11}+e_{31}^2/\epsilon_{33}$ is the answer for a plate constrained against thinning, which
is a different problem. Second, the thin-plate limit is a limit of the lowest branch only; the higher
branches do not go there.
\end{worked}

\subsection*{The bulk thermoelastic damping and the Hutson--White condition}

Two limits of the thermal side are worth having by hand, because both are used later as tests.

\subsection*{Fourier thermoelastic damping, in full}

\begin{worked}
Take a longitudinal plane wave in an unbounded medium: the carrier column is dropped, $\tau_q=0$ so that
$f=1$, and only $u_x$ and $\theta$ survive with the wave running along $x$ and no $z$ dependence, so
$\eta=0$ and $L=-k^2$.
Substitution: the momentum row is $(c_{11}k^2-\rho\omega^2)U+\ii k\beta_1\Theta=0$ and the heat row is
$T_0\beta_1\omega k\,U+(k_{\rm th}k^2-\ii\omega\rho c_p)\Theta=0$, the $\omega k$ in the second coming
from $\dot\varepsilon_{xx}=(-\ii\omega)(\ii k)U=+\omega kU$.
Collect: the determinant of that $2\times2$ system is
$(c_{11}k^2-\rho\omega^2)(k_{\rm th}k^2-\ii\omega\rho c_p)-\ii k\beta_1\,T_0\beta_1\omega k=0$.
Now perturb. Write $\omega=\omega_0(1-\ii\delta)$ with $\rho\omega_0^2=c_{11}k^2$ and $|\delta|\ll1$.
Then $c_{11}k^2-\rho\omega^2=\rho\omega_0^2-\rho\omega_0^2(1-\ii\delta)^2\approx2\ii\rho\omega_0^2\delta$.
For the heat factor, factor out $-\ii\omega\rho c_p$:
$k_{\rm th}k^2-\ii\omega\rho c_p=-\ii\omega\rho c_p\bigl(1+\ii k_{\rm th}k^2/(\omega\rho c_p)\bigr)
=-\ii\omega\rho c_p\bigl(1+\ii/(\omega\tau)\bigr)$ with
$\tau=\rho c_p/(k_{\rm th}k^2)$, the diffusion time of a thermal layer one wavelength thick.
Substitute both:
$2\ii\rho\omega_0^2\delta\,(-\ii\omega\rho c_p)(1+\ii/(\omega\tau))-\ii T_0\beta_1^2\omega k^2=0$,
that is $2\rho^2\omega_0^2\omega c_p\delta\,(1+\ii/(\omega\tau))=\ii T_0\beta_1^2\omega k^2$.
Collect: $\delta=\dfrac{\ii\,T_0\beta_1^2k^2}{2\rho^2\omega_0^2c_p(1+\ii/(\omega\tau))}
=\dfrac{\ii\Delta_E/2}{1+\ii/(\omega\tau)}$ with
$\Delta_E=\dfrac{T_0\beta_1^2k^2}{\rho c_p\,\rho\omega_0^2}=\dfrac{T_0\beta_1^2}{\rho c_p c_{11}}$,
using $\rho\omega_0^2=c_{11}k^2$.
Multiply above and below by $1-\ii/(\omega\tau)$ and take real and imaginary parts:
$\delta=\dfrac{\Delta_E}{2}\dfrac{\omega\tau+\ii(\omega\tau)^2}{1+(\omega\tau)^2}$, so
$\operatorname{Re}\delta=\dfrac{\Delta_E}{2}\dfrac{\omega\tau}{1+(\omega\tau)^2}$ and
$\operatorname{Im}\delta=\dfrac{\Delta_E}{2}\dfrac{(\omega\tau)^2}{1+(\omega\tau)^2}$.
Result: since $\omega=\omega_0(1-\ii\delta)$, the attenuation in time is
$-\operatorname{Im}\omega=\omega_0\operatorname{Re}\delta$ and the quality factor is
$Q^{-1}=2\operatorname{Re}\delta=\Delta_E\,\dfrac{\omega\tau}{1+(\omega\tau)^2}$,
with a velocity shift $\Delta c/c=\operatorname{Im}\delta=\dfrac{\Delta_E}{2}\dfrac{(\omega\tau)^2}{1+(\omega\tau)^2}$.
$Q^{-1}$ is largest at $\omega\tau=1$, that is at $\omega=\rho c_pc_{11}/k_{\rm th}$ once $k=\omega/c_L$
is used, and the peak value is $\Delta_E/2$. This is the standard Zener thermoelastic damping formula,
and the test is that the full matrix reproduces it in this limit.
\end{worked}

\subsection*{The Hutson--White condition, in full}

\begin{worked}
Keep the shear wave: $u_z$ with the field along $x$, carrier terms on, thermal terms off, no drift, no
diffusion, no recombination, and no $z$ dependence, so that $\eta=0$. The three rows that survive are the
$z$ momentum, Gauss' law and the carrier equation:
$(-\rho\omega^2+c_{44}k^2)U+e_{15}k^2\Phi=0$, $-e_{15}k^2U+\epsilon_{11}k^2\Phi+q\hat n=0$ and
$-\ii\omega\hat n-b_nn_0k^2\Phi=0$, the last from $\partial_t\Delta n=-b_nn_0\nabla^2\varphi$ with
$\nabla^2\to-k^2$.
Substitution: eliminate $\hat n$ from the third row, $\hat n=\ii b_nn_0k^2\Phi/\omega$, and put it into
the second: $-e_{15}U+\bigl(\epsilon_{11}+\ii qb_nn_0/\omega\bigr)k^2\Phi/k^2=0$, so
$\Phi=e_{15}U/\epsilon_{\rm eff}$ with
$\epsilon_{\rm eff}=\epsilon_{11}\bigl(1+\ii\omega_c/\omega\bigr)$ and
$\omega_c=qb_nn_0/\epsilon_{11}$.
Collect: the momentum row then gives $\rho\omega^2=k^2\bigl(c_{44}+e_{15}^2/\epsilon_{\rm eff}\bigr)$, so
$k=\omega\sqrt{\rho/(c_{44}+e_{15}^2/\epsilon_{\rm eff})}$, complex.
Result: write $K^2=e_{15}^2/(c_{44}\epsilon_{11})$ and expand for $K^2\ll1$:
$k\approx(\omega/c_{\rm sh})\bigl(1-\tfrac12 K^2/(1+\ii\omega_c/\omega)\bigr)$ with
$c_{\rm sh}=\sqrt{c_{44}/\rho}$, so the loss per wavelength is
$2\pi\operatorname{Im}k/\operatorname{Re}k\approx\pi K^2\dfrac{\omega/\omega_c}{1+(\omega/\omega_c)^2}$.
It peaks at $\omega=\omega_c$, that is at
$\dfrac{\epsilon_{11}\,\omega}{q\,b_n}=n_0$:
the Hutson--White condition. It locates the carrier density at which the attenuation is greatest for a
given frequency, and it is a property of the dielectric relaxation of the free charge, not of the plate.
At finite $K^2$ the peak moves by a relative amount of order $K^2$, which the weak-coupling statement does
not claim to cover.
\end{worked}

\section{Calculation G, continued: one column of $\mathbf B$ in full}

Each column of $\mathbf B$ belongs to one root, and each entry is one row function evaluated at that root
with the amplitudes of that root. Write the amplitude vector of root $m$ as
$X_m(1,\;r_z,\;r_\varphi,\;r_n,\;r_\theta)$, where the four ratios come from Calculation~E and the overall
scale is carried by $A_m$. Then every entry of the column is $X_m$ times an expression in $\eta_m$, $k$
and $\omega$ alone, and it is worth having those seven expressions on the page.

\begin{worked}
Substitution: put $Z_m=r_zX_m$, $\Phi_m=r_\varphi X_m$, $\nu_m=r_nX_m$, $\Theta_m=r_\theta X_m$ in the
seven rows of Section~\ref{sec:calcF}. Collect $X_m$ out of each.
$b_1=\Ncal_m^{-1}\bigl[\Gcal_m(\ii kc_{13}+\eta c_{33}r_z)+\eta e_{33}r_\varphi-\beta_3r_\theta
-l_g^2\ii k\eta\,c_{44}(\eta+\ii kr_z)\bigr]X_m$.
$b_2=\Ncal_m^{-1}\bigl[\Gcal_mc_{44}(\eta+\ii kr_z)+\ii ke_{15}r_\varphi
-l_g^2\ii k\eta(\ii kc_{11}+\eta c_{13}r_z)\bigr]X_m$.
Result: the two traction rows, each with its surface-gradient correction.
\end{worked}

\begin{worked}
$b_3=(l_g^2/\Ncal_m)\,\eta\,c_{44}(\eta+\ii kr_z)X_m$ and
$b_4=(l_g^2/\Ncal_m)\,\eta\,(\ii kc_{13}+\eta c_{33}r_z)X_m$. Both carry $l_g^2$ and both vanish with it,
which is the statement that the two rows drop out of the count when the gradient length is zero.
\end{worked}

\begin{worked}
$b_5=\bigl[\ii ke_{31}+\eta e_{33}r_z-\eta\epsilon_{33}r_\varphi+p_3r_\theta\bigr]X_m$, the open-circuit
row. $b_6=\eta\bigl[-D_nr_n+b_nn_0r_\varphi-\beta_{\rm th}r_\theta\bigr]X_m$, the zero-flux row; the factor
$\eta$ in front is worth noticing, because it means the row vanishes at $\eta=0$.
$b_7=\eta\,r_\theta X_m$, the adiabatic row. Result: the last three rows.
\end{worked}

\begin{worked}
Now assemble. The column of $\mathbf B$ belonging to root $m$ is
$\bigl(b_1\ldots b_7\bigr)^{T}e^{\eta_mh}$ in the upper block and
$\bigl(b_1\ldots b_7\bigr)^{T}e^{-\eta_mh}$ in the lower one, with the same seven expressions. Multiplying
a column by a non-zero number multiplies $\det\mathbf B$ by that number, so the common factor $X_m$ may be
dropped from every column at once, and the column scaling of Section~\ref{sec:calcG} is the same operation
applied with $e^{-|\operatorname{Re}\eta_m|h}$. Result: a matrix whose entries are the seven expressions
above, evaluated at the fourteen roots, times the two exponentials.
\end{worked}

\section{Calculation I: the dimension of every entry}

The check costs ten minutes and catches the commonest error in a reduction of this length, which is a
factor of $\omega$ or of $\ii$ lost while moving a term across. Every row of $M(\eta)\mathbf a=\mathbf 0$
is one physical balance, so every entry in that row must have the dimension that turns its amplitude into
the dimension of the balance.

\begin{worked}
Row 1 is a force balance per unit volume, N\,m$^{-3}$. Check each entry.
$M_{11}X$: $\Gcal c_{11}k^2$ is Pa\,m$^{-2}$ = N\,m$^{-4}$, and $X$ is m, so the product is N\,m$^{-3}$.
$M_{12}Z$: $\ii k\eta\Gcal(c_{13}+c_{44})$ is m$^{-2}\,$Pa, and $Z$ is m, so again N\,m$^{-3}$.
$M_{13}\Phi$: $\ii k\eta(e_{31}+e_{15})$ is m$^{-2}$\,C\,m$^{-2}$ = C\,m$^{-4}$, and $\Phi$ is V, so the
product is C\,V\,m$^{-4}$ = J\,m$^{-4}$ = N\,m$^{-3}$.
$M_{15}\Theta$: $\ii k\beta_1$ is m$^{-1}$\,Pa\,K$^{-1}$, and $\Theta$ is K, so the product is
Pa\,m$^{-1}$ = N\,m$^{-3}$. Row 1 is consistent.
\end{worked}

\begin{worked}
Row 3 is Gauss' law, C\,m$^{-3}$. $M_{31}X$: $\ii k\eta(e_{15}+e_{31})$ is C\,m$^{-4}$ and $X$ is m, so
C\,m$^{-3}$. $M_{32}Z$: $e_{33}\eta^2$ is C\,m$^{-2}$\,m$^{-2}$ = C\,m$^{-4}$ and $Z$ is m, so C\,m$^{-3}$.
$M_{33}\Phi$: $\epsilon_{11}k^2$ is C\,V$^{-1}$m$^{-1}$\,m$^{-2}$ and $\Phi$ is V, so C\,m$^{-3}$.
$M_{34}\hat n$: $q$ is C and $\hat n$ is m$^{-3}$, so C\,m$^{-3}$.
$M_{35}\Theta$: $\eta p_3$ is m$^{-1}$\,C\,m$^{-2}$K$^{-1}$ and $\Theta$ is K, so C\,m$^{-3}$. Row 3 is
consistent.
\end{worked}

\begin{worked}
Row 4 is a number balance per unit volume per unit time, m$^{-3}$s$^{-1}$.
$M_{43}\Phi$: $b_nn_0L$ is m$^2$V$^{-1}$s$^{-1}\cdot$ m$^{-3}\cdot$ m$^{-2}$ = V$^{-1}$m$^{-3}$s$^{-1}$ and
$\Phi$ is V, so m$^{-3}$s$^{-1}$.
$M_{44}\hat n$: every term in $M_{44}$ is a rate, s$^{-1}$, and $\hat n$ is m$^{-3}$.
$M_{45}\Theta$: $\beta_{\rm th}L$ needs $\beta_{\rm th}$ in m\,s$^{-1}$K$^{-1}$, which is the unit that
makes $\beta_{\rm th}\nabla^2\theta$ a carrier flux divergence; times $\Theta$ in K it gives
m$^{-3}$s$^{-1}$. Row 4 is consistent.
\end{worked}

\begin{worked}
Row 5 is an energy balance per unit volume per unit time, J\,m$^{-3}$s$^{-1}$.
$M_{51}X$: $fT_0\omega k\beta_1$ is K$\,\cdot$s$^{-1}\cdot$m$^{-1}\cdot$Pa\,K$^{-1}$ = Pa\,m$^{-1}$s$^{-1}$,
and $X$ is m, so Pa\,s$^{-1}$ = J\,m$^{-3}$s$^{-1}$.
$M_{52}Z$: $-\ii\omega fT_0\beta_3\eta$ is s$^{-1}\,$K$\,$Pa\,K$^{-1}$m$^{-1}$ = Pa\,m$^{-1}$s$^{-1}$, and
$Z$ is m, so the same.
$M_{53}\Phi$: $\ii\omega fT_0p_3\eta$ is
s$^{-1}$K$\,$C\,m$^{-2}$K$^{-1}$m$^{-1}$ = C\,m$^{-3}$s$^{-1}$ and $\Phi$ is V, so J\,m$^{-3}$s$^{-1}$.
$M_{54}\hat n$: the Peltier term $-\ii\omega q\Pi_{\rm th}$ is
s$^{-1}\,$C$\,\cdot$V, so $\Pi_{\rm th}$ carries the unit of a thermoelectric power, V; the recombination
term $E_g/\tau_R$ is an energy per recombination, J, divided by a time, s. Both are J\,s$^{-1}$ per carrier,
and multiplied by $\hat n$ in m$^{-3}$ both give J\,m$^{-3}$s$^{-1}$.
$M_{55}\Theta$: $-\ii\omega f\rho c_p$ is s$^{-1}$J\,m$^{-3}$K$^{-1}$ and $-k_{\rm th}L$ is
W\,m$^{-1}$K$^{-1}$m$^{-2}$ = W\,m$^{-3}$K$^{-1}$; times $\Theta$ in K both give J\,m$^{-3}$s$^{-1}$. Row 5
is consistent.
\end{worked}

\section{Calculation J: what the failure of the split costs}

\begin{worked}
When the three thermal--carrier couplings vanish, the fourteen roots pair as $\pm\eta_m$ and the fourteen
conditions split into two sets of seven. Write the amplitudes of the pair as related by the column signs
$\tau=(+1,-1,-1,-1,+1)$ of Section~\ref{sec:calcD}: $\mathbf a(-\eta_m)=\operatorname{diag}(\tau)\,
\mathbf a(\eta_m)$. Substitution: form the two combinations
$\mathbf a(\eta_m)e^{\eta_mz}\pm\operatorname{diag}(\tau)\mathbf a(\eta_m)e^{-\eta_mz}$. Collect: component
$j$ of the plus combination is $2a_j\cosh\eta_mz$ when $\tau_j=+1$ and $2a_j\sinh\eta_mz$ when
$\tau_j=-1$. Result: a solution of definite parity, for which the conditions at $z=-h$ follow from those at
$z=+h$, so seven conditions at one face and seven roots give a $7\times7$ matrix per family.
\end{worked}

\begin{worked}
When $\beta_{\rm th}$, $\Pi_{\rm th}$ or $E_g$ is present, the parity argument fails and there is no such
combination. Substitution: nothing to substitute. Collect: fourteen roots, fourteen conditions, one
$14\times14$ matrix, and no saving. Result: the same dispersion relation written the long way, and the
numerical work has to carry fourteen columns instead of seven.
\end{worked}

The practical consequences are three, and they belong in the paper. First, the two routes must agree when
the three couplings are set to zero, and checking that is one of the self-tests: run the fourteen-by-fourteen
and the two seven-by-seven at the same constants and compare the first three branches. Second, the
fourteen-by-fourteen is the one to use for the figures, because the coupling is the point of the paper.
Third, the ordering of the roots matters for the conditioning and not for the physics: sort them by
$|\operatorname{Re}\eta_m|$ before assembling the matrix, and scale each column by
$e^{-|\operatorname{Re}\eta_m|h}$, which multiplies the determinant by a non-zero number and leaves its
zeros where they are.

\begin{worked}
One more remark on the numerical side, because it is a fact about the matrix and not an algorithm. The
thermal and carrier roots are large: with the constants of the internal test set, the thermal pair sits
near $10^{8}$\,m$^{-1}$ while the elastic pair sits near $10^{4}$\,m$^{-1}$. Substitution into
$e^{\eta_mh}$ at a millimetre thickness gives $e^{10^{5}}$, and into $e^{-\eta_mh}$ gives its reciprocal,
so the raw entries of $\mathbf B$ span $10^{5}$ decades and no floating-point determinant survives them.
Collect: the column scaling above is what makes the determinant computable at all. Result: the zeros of
$\det\mathbf B$ are where they always were; only the arithmetic becomes possible.
\end{worked}

\subsection*{Why each test is a test}

\begin{worked}
The elastic test is the only one with an independent answer, so it comes first and everything else is
checked against our own algebra. The parity test is decided by the parities of the entries alone, so it is
a test of the reduction and not of the physics: if the three thermal--carrier couplings are set to zero and
the roots do not pair, the error is in one entry, and the entries to look at are the ones that carry
$\beta_{\rm th}$, $\Pi_{\rm th}$ and $E_g$. The open-circuit and shorted thin-plate speeds bracket the
coupling and are both measurable, so they test the piezo entries against a laboratory number. The Zener
curve tests the thermal entries twice over: the shape $\omega\tau/(1+\omega^2\tau^2)$ tests the coupling
and the position of the peak tests the relaxation time. The Hutson--White condition tests the carrier
entries, and because it is a condition on the carrier density it can be checked by sweeping $n_0$, which is
a knob no other test turns.
Result: each test fails in a different place, and the order above is the order in which they should be run.
\end{worked}

\section{Calculation K: why seven conditions and not five}

\begin{worked}
Count the amplitudes first. The degree of $\det M$ is the number of partial waves, because each root
$\eta_m$ contributes one wave and one unknown amplitude $A_m$. With $l_g\ne0$ the count of
Section~\ref{sec:calcD} gives fourteen, with $l_g=0$ it gives ten. Each face supplies half of them, so
seven conditions per face in the first case and five in the second.
Now count the other way, by field, and check that the two agree. There are five fields, so five conditions
per face is one per field: a traction, a shear traction, an electric condition, a carrier or phason
condition and a thermal condition. The two extra conditions in the gradient case are the two double
tractions, and they are extra because the gradient length raises the order of the mechanical equations in
$z$ from two to four: $\Gcal$ carries $\eta^2$ and multiplies the $\eta^2$ that the strains already carry,
so $M_{11}$ and $M_{22}$ are quartic in $\eta$ and the two momentum rows are fourth order in $z$.
Collect: five fields plus two extra orders is seven, and seven per face times two faces is fourteen, which
is the number of amplitudes. The two counts agree, which is the check.
\end{worked}

\begin{worked}
The same count says what happens at the degenerate edges. At $l_g=0$ the double tractions vanish and the
matrix drops to ten, so five conditions per face; keeping the two double-traction rows with $l_g=0$ would
give two rows of zeros and a singular matrix, which is the commonest way to get a spurious zero branch in
this problem. At $\eta=0$ the carrier flux row vanishes identically, because every term in it carries one
$\eta$, so a root at the origin has to be treated on its own rather than through the general assembly.
Result: the count is not a formality; it is where the degenerate cases announce themselves.
\end{worked}

\section{Self-test list}\label{sec:selftest}

None of these needs a number from me.

\begin{enumerate}
\item Set every coupling to zero and both size lengths to zero: the dispersion must be the classical
Rayleigh--Lamb one for an isotropic plate, both families, and the lowest antisymmetric branch must tend to
the Rayleigh speed at large $f(2h)$.
\item Keep only the piezoelectric coupling: the thin-plate limit must give the open-circuit stiffened bulk
speed of \eqref{eq:heat}, and the shorted limit the unstiffened one.
\item Keep the heat equation in Fourier form, $\tau_q=0$, and switch the carrier off: a longitudinal bulk
wave must show thermoelastic damping with the standard dependence on $\omega\tau$.
\item Keep only the carrier terms with the piezoelectric shear coupling: the bulk attenuation must peak at
the carrier density predicted by the relaxation of the free charge.
\item Let $l_g$ and $l_{\rm nl}$ tend to zero separately: the roots must approach the classical ones
continuously, the gradient length must stiffen, and the nonlocal length must soften.
\item Count the roots of $\det M$ for one case with $l_g\ne0$ and one with $l_g=0$ and compare with the
number of rows you use in $\mathbf B$. A mismatch is a counting error, not a physical finding.
\item Switch $\beta_{\rm th}$, $\Pi_{\rm th}$ and $E_g$ off together and check that the roots come in
pairs $\pm\eta_m$. Switch any one of the three back on and check that the pairing is lost. Then switch
$p_3$ alone, on and off, and check that the pairing is unaffected, which is the statement of
Section~\ref{sec:calcD}.
\item Check the dimensions of every entry of row 5 against row 1. A factor of $\omega$ is easy to lose
in the Moore--Gibson--Thompson factor.
\end{enumerate}

\section{Work plan}\label{sec:plan}

Weeks one to two: read the sources marked in Section~\ref{sec:source}, in particular the two equations of
Lim et al.\ and the drift--diffusion pair of Yang and Zhou, and write the entry-by-entry matrix on paper
from the field equations alone. Weeks three to four: the roots, the amplitude ratios, the parity check for
both settings of $p_3$, and the counting of rows against roots. Weeks five to six: the limits of
Calculation H, in the order given, and the comparison of the elastic limit with the classical plate.
Weeks seven to eight: the dispersion curves against $f(2h)$, the attenuation of the first three branches,
and the thickness map that shows where the relaxation time is resolvable. Weeks nine to ten: the size
effect, one figure for the two lengths. Weeks eleven to twelve: writing, figures, and the checking of every
citation against the page it came from.

\bibliographystyle{unsrt}
\bibliography{Lamb_Sumit}

\end{document}
